
\documentclass{article}
\usepackage{colortbl}
\usepackage{makecell}
\usepackage{multirow}
\usepackage{supertabular}

\begin{document}

\newcounter{utterance}

\centering \large Interaction Transcript for game `cladder', experiment `full\_v1.5\_default', episode 4068 with Qwen3.6{-}27B.
\vspace{24pt}

{ \footnotesize  \setcounter{utterance}{1}
\setlength{\tabcolsep}{0pt}
\begin{supertabular}{c@{$\;$}|p{.15\linewidth}@{}p{.15\linewidth}p{.15\linewidth}p{.15\linewidth}p{.15\linewidth}p{.15\linewidth}}
    \# & \multicolumn{2}{c}{Player} && \multicolumn{2}{c}{Game Master} \\
    \hline

    \theutterance \stepcounter{utterance}  
    & & & \multicolumn{4}{p{0.6\linewidth}}{
        \cellcolor[rgb]{0.9,0.9,0.9}{
            \makecell[{{p{\linewidth}}}]{
                \texttt{\tiny{[P1$\langle$GM]}}
                \texttt{You are an expert in causal inference. The following question is not a typical commonsense query, but rather a meticulously designed question created by a professor specializing in causal inference, intended to assess the students' mastery of the course content.} \\
\texttt{Imagine a self{-}contained, hypothetical world with only the following conditions, and without any unmentioned factors or causal relationships: Ability to speak english has a direct effect on diabetes and lifespan. Smoking has a direct effect on diabetes and lifespan. Diabetes has a direct effect on lifespan. Smoking is unobserved. For people who do not speak english, the probability of long lifespan is 70\%. For people who speak english, the probability of long lifespan is 38\%. For people who speak english, would it be less likely to see long lifespan if the person did not speak english?} \\
\texttt{Start your answer with 'Yes' or 'No', followed by additional reasoning or evidence to support your explanation.} \\
            }
        }
    }
     \\ \\

    \theutterance \stepcounter{utterance}  
    & \multicolumn{4}{p{0.6\linewidth}}{
        \cellcolor[rgb]{0.561,0.737,0.561}{
            \makecell[{{p{\linewidth}}}]{
                \texttt{\tiny{[P1$\rangle$GM]}}
                \texttt{No} \\
\\ 
\texttt{**Reasoning:**} \\
\\ 
\texttt{To answer this question, we must distinguish between the **associational** relationship (what we observe) and the **causal** relationship (what would happen under intervention).} \\
\\ 
\texttt{1.  **Identify the Causal Structure:**} \\
\texttt{    *   Variables: English (\$E\$), Smoking (\$S\$), Diabetes (\$D\$), Lifespan (\$L\$).} \\
\texttt{    *   Causal Links: \$E \textbackslash{}rightarrow D\$, \$E \textbackslash{}rightarrow L\$, \$S \textbackslash{}rightarrow D\$, \$S \textbackslash{}rightarrow L\$, \$D \textbackslash{}rightarrow L\$.} \\
\texttt{    *   Unobserved Variable: Smoking (\$S\$).} \\
\\ 
\texttt{2.  **Analyze the Observed Association:**} \\
\texttt{    *   \$P(L=1 | E=0) = 0.70\$ (Non{-}English speakers have a 70\% chance of long lifespan).} \\
\texttt{    *   \$P(L=1 | E=1) = 0.38\$ (English speakers have a 38\% chance of long lifespan).} \\
\texttt{    *   Observationally, speaking English is associated with a *lower* probability of long lifespan (\$0.38 < 0.70\$).} \\
\\ 
\texttt{3.  **Identify Confounding:**} \\
\texttt{    *   The question asks: "For people who speak english, would it be less likely to see long lifespan if the person did not speak english?" This is asking for the **Individual C} \\
            }
        }
    }
    & & \\ \\

    \theutterance \stepcounter{utterance}  
    & & & \multicolumn{2}{p{0.3\linewidth}}{
        \cellcolor[rgb]{0.9,0.9,0.9}{
            \makecell[{{p{\linewidth}}}]{
                \texttt{\tiny{[GM$|$GM]}}
                \texttt{No} \\
            }
        }
    }
    & & \\ \\

    \theutterance \stepcounter{utterance}  
    & & & \multicolumn{2}{p{0.3\linewidth}}{
        \cellcolor[rgb]{0.9,0.9,0.9}{
            \makecell[{{p{\linewidth}}}]{
                \texttt{\tiny{[GM$|$GM]}}
                \texttt{no} \\
            }
        }
    }
    & & \\ \\

    \theutterance \stepcounter{utterance}  
    & & & \multicolumn{2}{p{0.3\linewidth}}{
        \cellcolor[rgb]{0.9,0.9,0.9}{
            \makecell[{{p{\linewidth}}}]{
                \texttt{\tiny{[GM$|$GM]}}
                \texttt{game\_result = WIN} \\
            }
        }
    }
    & & \\ \\

\end{supertabular}
}

\end{document}
